Un fil infinit que porta un corrent de 2~A es troba a 5,0~cm de distància del centre d'una espira circular de 2,0~cm de diàmetre que transporta 500~mA.

\figura[0.45]{fig1.png}

\begin{apartat}{1}
Calculeu el vector del camp magnètic al centre de l'espira produït pel fil infinit i el vector del camp magnètic al centre de l'espira que produeix la mateixa espira.
\end{apartat}

\begin{apartat}{1}
Quin és el valor del camp magnètic total al centre de l'espira? Si volem un camp magnètic total $B = 0$ al centre de l'espira, quin ha de ser el valor de la nova intensitat que hi circuli?
\end{apartat}

\blocetiquetat{Dada:}{$\mu_0 = 4\pi\times 10^{-7}\ \mathrm{T\,m\,A^{-1}}$}
\nota{El mòdul del camp magnètic creat per un fil infinit pel qual circula una intensitat $I$ és: $B = \dfrac{\mu_0 I}{2\pi r}$, on $r$ és la distància al fil conductor. El mòdul del camp magnètic al centre d'una espira de corrent de radi $R$ és: $B = \dfrac{\mu_0 I}{2R}$.}
